\chapter*{Use of generative AI}
\addcontentsline{toc}{chapter}{Use of generative AI}
\label{ch:UseOfGenerativeAI}

Claude Code by Anthropic assisted in the development of two components of the codebase utilized in this thesis.
The first component, developed with the assistance of the Claude Opus 4.8 model, manages operations involving jar files.
The \texttt{spec} package within the migration process opens the jar file of a compiled rule set and instantiates the contained change propagation specifications \cite{vitruv-dsls-spec-package}.
\texttt{ReactionsServerLauncher} serves as the primary class of the language server jar, while \texttt{ReactionsLanguageIdeSetup} registers the metamodels included in the jar \cite{vitruv-dsls-ide-package}.
The build process for the Visual Studio Code extension packages this jar, and the extension launches it alongside additional metamodel jars \cite{vitruv-dsls-vscode-plugin}.
The second component, developed with the assistance of the Claude Opus 5 model, consists of the scripts located in the \texttt{utils} folder of \cite{thesis-repo}, as listed in \autoref{table:evaluation-scripts}.
This component also includes \texttt{generate-config-matrix}, \texttt{generate-selection-matrix}, and the scripts in the \texttt{reversibility} folder of \cite{vitruv-dsls-fork}, which generate the evaluated data.
All generated code underwent review prior to its utilization.
